\chapter{Estado del Arte}

\section{El Modelo de Media-Varianza de Markowitz}
En 1952, Harry Markowitz introdujo el modelo de media-varianza para la optimización de carteras En el modelo de \cite{markowitz1952}, el autor toma la varianza como la media del riesgo. Definiendo cada activo como una variable aleatoria, teniendo en cosndieración el retorno esperado para un tiempo determinado, así como la variabilidad medida a través de la varianza de dicho retorno. De forma pionera, se desmiente la idea de que el inversor deba limitarse a buscar el máximo retorno esperado a la hora de confeccionar su cartera. Si uno se guía únicamente por esta premisa, Markowtiz nos muestra que no tiene en consideración el riesgo asociado a los activos seleccionados, así como la correlación entre los retornos de los mismos. 


Es por ello que el artículo se centra en la optimización de carteras buscando aquellas combinaciones de activos que, para igual retorno esperado, minimicen la varianza, es decir, el riesgo. Por lo tanto, el problema de optimizar una cartera se reduce a encontrar máximos o mínimos dentro de una función de dos variables, fijando una de ellas. Teniendo en cuenta solo aquellas combinaciones de activos que se puedan realizar, las combinaciones deseadas forman lo que el autor denomina la frontera eficiente. 



\section{Limitaciones de la Teoría Clásica}
Markowitz fue pionero, sin embargo su modelo presenta limitaciones aparentes, que con las herramientas actuales de análisis de datos, se pueden superar. 

El modelo asume que los retornos de los activos financieros siguen una distribución normal. Sin embargo, la evidencia empírica muestra que no siempre se puede asumir la normalidad en los retornos de los activos. Por otro lado, las correlaciones entre activos financieros no son estáticas, sino que varían a lo largo del tiempo.

\section{Fundamentos de la Teoría de la Información}
